\subsection{MSVC}

\RU{Компилируем в}\EN{Let's compile it in}\ES{Compilémoslo en} MSVC 2010:

\begin{lstlisting}
cl 1.cpp /Fa1.asm
\end{lstlisting}

\RU{(Ключ /Fa означает сгенерировать листинг на ассемблере)}
\EN{(/Fa option instructs the compiler to generate assembly listing file)}\ES{(La opción /Fa le indica al compilador que genere un archivo de listado en ensamblador)}

\begin{lstlisting}[caption=MSVC 2010]
CONST	SEGMENT
$SG3830	DB	'hello, world', 0AH, 00H
CONST	ENDS
PUBLIC	_main
EXTRN	_printf:PROC
; Function compile flags: /Odtp
_TEXT	SEGMENT
_main	PROC
	push	ebp
	mov	ebp, esp
	push	OFFSET $SG3830
	call	_printf
	add	esp, 4
	xor	eax, eax
	pop	ebp
	ret	0
_main	ENDP
_TEXT	ENDS
\end{lstlisting}

\ifx\LITE\undefined
\RU{MSVC выдает листинги в синтаксисе Intel.}\EN{MSVC produces assembly listings in Intel-syntax.}\ES{MSVC genera los listados de ensamblador en sintaxis Intel.} 
\RU{Разница между синтаксисом Intel и AT\&T будет рассмотрена немного позже:}
\EN{The difference between Intel-syntax and AT\&T-syntax will be discussed in}\ES{La diferencia entre la sintaxis Intel y la sintaxis AT\&T se discutirá en} \myref{ATT_syntax}.
\fi

\RU{Компилятор сгенерировал файл \TT{1.obj}, который впоследствии будет слинкован линкером в \TT{1.exe}.}
\EN{The compiler generated \TT{1.obj} file, which is to be linked into \TT{1.exe}.}\ES{El compilador generó el archivo \TT{1.obj}, el cual será enlazado en \TT{1.exe}.}
\RU{В нашем случае этот файл состоит из двух сегментов: \TT{CONST} (для данных-констант) и \TT{\_TEXT} (для кода).}
\EN{In our case, the file contains two segments: \TT{CONST} (for data constants) and \TT{\_TEXT} (for code).}\ES{En nuestro caso, el archivo contiene dos segmentos: \TT{CONST} (para constantes de datos) y \TT{\_TEXT} (para código).}

\index{\CLanguageElements!const}
\label{string_is_const_char}
\RU{Строка \TT{hello, world} в \CCpp имеет тип \TT{const char[]} \cite[p176, 7.3.2]{TCPPPL}, 
однако не имеет имени.}
\EN{The string \TT{hello, world} in \CCpp has type \TT{const char[]} \cite[p176, 7.3.2]{TCPPPL},
however it does not have its own name.}\ES{La cadena \TT{hello, world} en \CCpp tiene el tipo \TT{const char[]} \cite[p176, 7.3.2]{TCPPPL}; sin embargo, no tiene un nombre propio.}
\RU{Но компилятору нужно как-то с ней работать, поэтому он дает ей внутреннее имя \TT{\$SG3830}.}%
\EN{The compiler needs to deal with the string somehow so it defines the internal name \TT{\$SG3830} for it.}\ES{El compilador necesita manejar la cadena de alguna manera, por lo que le define el nombre interno \TT{\$SG3830}.}%

\RU{Поэтому пример можно было бы переписать вот так}\EN{That is why the example may be rewritten as follows}\ES{Por eso el ejemplo podría reescribirse de la siguiente manera}:

\lstinputlisting{patterns/01_helloworld/hw_2.c}

\RU{Вернемся к листингу на ассемблере. Как видно, строка заканчивается нулевым байтом~--- 
это требования стандарта \CCpp для строк}%
\EN{Let's go back to the assembly listing. 
As we can see, the string is terminated by a zero byte, which is standard for \CCpp strings}.\ES{Volvamos al listado en ensamblador. Como podemos ver, la cadena termina con un byte cero, lo cual es estándar para las cadenas en \CCpp}.
\RU{Больше о строках в Си}\EN{More about C strings}: \myref{C_strings}.\ES{Más sobre las cadenas en C: \myref{C_strings}.}

\RU{В сегменте кода \TT{\_TEXT} находится пока только одна функция}%
\EN{In the code segment, \TT{\_TEXT}, there is only one function so far}: \main.\ES{En el segmento de código, \TT{\_TEXT}, solo hay una función por ahora: \main.}
\RU{Функция \main, как и практически все функции, начинается с пролога и заканчивается эпилогом}%
\EN{The function \main starts with prologue code and ends with epilogue code (like almost any function)}%
\footnote{\RU{Об этом смотрите подробнее в разделе о прологе и эпилоге функции}%
\EN{You can read more about it in section about function prolog and epilog}%
~(\myref{sec:prologepilog}).}.\ES{La función \main comienza con el código de prólogo y termina con el de epílogo (como casi cualquier función)}\footnote{\ES{Puedes leer más sobre esto en la sección sobre el prólogo y epílogo de la función}~(\myref{sec:prologepilog}).}.

\index{x86!\Instructions!CALL}
\RU{Далее следует вызов функции \printf}
\EN{After the function prologue we see the call to the \printf function}: \TT{CALL \_printf}.\ES{Después del prólogo de la función vemos la llamada a la función \printf: \TT{CALL \_printf}.}
\index{x86!\Instructions!PUSH}
\RU{Перед этим вызовом адрес строки (или указатель на неё) с нашим приветствием при помощи инструкции \PUSH помещается в стек.}
\EN{Before the call the string address (or a pointer to it) containing our greeting is placed on the stack with the help of the \PUSH instruction.}\ES{Antes de la llamada, la dirección de la cadena (o un puntero a ella) que contiene nuestro saludo se coloca en la pila con la ayuda de la instrucción \PUSH.}

\RU{После того, как функция \printf возвращает управление в функцию \main, адрес строки (или указатель на неё) всё ещё лежит в стеке.}
\EN{When the \printf function returns the control to the \main function, the string address (or a pointer to it) is still on the stack.}\ES{Cuando la función \printf devuelve el control a la función \main, la dirección de la cadena (o un puntero a ella) todavía está en la pila.}
\RU{Так как он больше не нужен, то \glslink{stack pointer}{указатель стека} (регистр \ESP) корректируется.} 
\EN{Since we do not need it anymore, the \gls{stack pointer} (the \ESP register) needs to be corrected.}\ES{Como ya no lo necesitamos, el \gls{stack pointer} (el registro \ESP) debe ser corregido.}

\index{x86!\Instructions!ADD}
\TT{ADD ESP, 4} \RU{означает прибавить 4 к значению в регистре \ESP.}
\EN{means add 4 to the \ESP register value.}\ES{significa sumar 4 al valor del registro \ESP.}
\RU{Почему 4? Так как это 32-битный код, для передачи адреса нужно 4 байта. В x64-коде это 8 байт.}
\EN{Why 4? Since this is 32-bit program, we need exactly 4 bytes for address passing through the stack. 
If it was x64 code we would need 8 bytes.}\ES{¿Por qué 4? Como este es un programa de 32 bits, necesitamos exactamente 4 bytes para pasar la dirección a través de la pila. Si fuera código x64, necesitaríamos 8 bytes.}
\TT{ADD ESP, 4} \RU{эквивалентно \TT{POP регистр}, но без использования какого-либо регистра\footnote{Флаги
процессора, впрочем, модифицируются}.}
\EN{is effectively equivalent to \TT{POP register} but without using any register\footnote{CPU flags, however, are modified}.}\ES{es efectivamente equivalente a \TT{POP registro} pero sin usar ningún registro\footnote{Sin embargo, las banderas de la CPU se modifican}.}

\index{Intel C++}
\index{\oracle}
\index{x86!\Instructions!POP}

\RU{Некоторые компиляторы, например, Intel C++ Compiler, в этой же ситуации могут вместо 
\ADD сгенерировать \TT{POP ECX} (подобное можно встретить, например, в коде \oracle{}, им скомпилированном),
что почти то же самое, только портится значение в регистре \ECX.}
\EN{Some compilers (like the Intel C++ Compiler) for the same purpose may emit \TT{POP ECX} 
instead of \ADD (e.g., such a pattern can be observed in the \oracle{} code as it is compiled with the Intel C++ compiler).
This instruction has almost the same effect but the \ECX register contents will be overwritten.}\ES{Algunos compiladores (como Intel C++ Compiler) para el mismo propósito pueden emitir \TT{POP ECX} en lugar de \ADD (por ejemplo, tal patrón se puede observar en el código de \oracle{} compilado con el compilador Intel C++). Esta instrucción tiene casi el mismo efecto, pero el contenido del registro \ECX se sobrescribirá.}
\RU{Возможно, компилятор применяет \TT{POP ECX}, потому что эта инструкция короче (1 байт у \TT{POP} против 3 у \TT{ADD}).}
\EN{The Intel C++ compiler probably uses \TT{POP ECX} since this instruction's opcode is shorter than 
\TT{ADD ESP, x} (1 byte for \TT{POP} against 3 for \TT{ADD}).}\ES{El compilador Intel C++ probablemente usa \TT{POP ECX} ya que el opcode de esta instrucción es más corto que \TT{ADD ESP, x} (1 byte para \TT{POP} frente a 3 para \TT{ADD}).}

\RU{Вот пример использования \TT{POP} вместо \TT{ADD} из \oracle}\EN{Here is an example of using \TT{POP} instead of \TT{ADD} from \oracle}\ES{Aquí hay un ejemplo de uso de \TT{POP} en lugar de \TT{ADD} en \oracle}:

\begin{lstlisting}[caption=\oracle 10.2 Linux (\RU{файл }app.o\EN{ file})]
.text:0800029A                 push    ebx
.text:0800029B                 call    qksfroChild
.text:080002A0                 pop     ecx
\end{lstlisting}

%\RU{О стеке можно прочитать в соответствующем разделе}
%\EN{Read more about the stack in section}~(\myref{sec:stack}).
\index{\CLanguageElements!return}
\RU{После вызова \printf в оригинальном коде на \CCpp указано \TT{return 0}~--- вернуть 0 
в качестве результата функции \main.} 
\EN{After calling \printf, the original \CCpp code contains the statement \TT{return 0}~---return 0 as the result of the \main function.}\ES{Después de llamar a \printf, el código original en \CCpp contiene la instrucción \TT{return 0}~---devolver 0 como resultado de la función \main.}
\index{x86!\Instructions!XOR}
\RU{В сгенерированном коде это обеспечивается инструкцией}
\EN{In the generated code this is implemented by the instruction}\ES{En el código generado esto se implementa mediante la instrucción} \TT{XOR EAX, EAX}.
\index{x86!\Instructions!MOV}
\RU{\XOR, как легко догадаться~--- \q{исключающее ИЛИ}}%
\EN{\XOR is in fact just \q{eXclusive OR}}\ES{\XOR es en realidad solo un \q{O exclusivo}}%
\footnote{\href{http://go.yurichev.com/17118}{wikipedia}}
\RU{, но компиляторы часто используют его вместо простого}
\EN{but the compilers often use it instead of}\ES{, pero los compiladores a menudo lo usan en lugar de}
\TT{MOV EAX, 0}\EMDASH{}\RU{снова потому, что опкод короче (2 байта у \TT{XOR} против 5 у \TT{MOV}).}
\EN{again because it is a slightly shorter opcode (2 bytes for \TT{XOR} against 5 for \TT{MOV}).}\ES{nuevamente porque es un opcode ligeramente más corto (2 bytes para \TT{XOR} frente a 5 para \TT{MOV}).}

\index{x86!\Instructions!SUB}
\RU{Некоторые компиляторы генерируют}\EN{Some compilers emit}\ES{Algunos compiladores emiten}
\INS{SUB EAX, EAX},
\RU{что значит \IT{отнять значение в} \EAX \IT{от значения в }\EAX,
что в любом случае даст 0 в результате.}
\EN{which means \IT{SUBtract the value in the} \EAX \IT{from the value in} \EAX,
which in any case resulting in zero.}\ES{lo que significa \IT{restar el valor de} \EAX \IT{del valor de }\EAX, lo que en cualquier caso da como resultado cero.}

\index{x86!\Instructions!RET}
\RU{Самая последняя инструкция \RET возвращает управление в вызывающую функцию.
Обычно это код \CCpp \ac{CRT}, который, в свою очередь, 
вернёт управление операционной системе.}
\EN{The last instruction \RET returns the control to the \gls{caller}.
Usually, this is \CCpp \ac{CRT} code, which in its turn returns the control to the \ac{OS}.}\ES{La última instrucción \RET devuelve el control a la \gls{caller}. Usualmente, este es el código de la \ac{CRT} de \CCpp, el cual a su vez devuelve el control al \ac{OS}.}
